\clearpage
\VerdierPhysicalPageAnchor{0206}%
\begingroup\small\setlength{\parskip}{0pt}\setlength{\emergencystretch}{1.5em}
\setlength{\abovedisplayskip}{4pt plus 2pt minus 2pt}%
\setlength{\belowdisplayskip}{4pt plus 2pt minus 2pt}%

\VerdierSemanticAnchor{ch3d-2-4-5}
\noindent\textbf{注意 2.4.5.}\quad
a) $\mathcal A$ を、そのすべての対象が射影対象の商対象と同型であるアーベル圏とする。
射影型複体は必ずしも $Q$-左自由ではない。

\VerdierSemanticAnchor{ch3d-2-4-5-example}
\emph{例}：$\mathcal A$ を $\mathbb Z/4\mathbb Z$-加群の圏とする。複体
\[
\cdots\longrightarrow\mathbb Z/4\mathbb Z
\longrightarrow\mathbb Z/4\mathbb Z
\longrightarrow\mathbb Z/4\mathbb Z
\longrightarrow\mathbb Z/4\mathbb Z
\longrightarrow\mathbb Z/4\mathbb Z\longrightarrow\cdots
\]
のすべての微分が $2$ 倍写像であるとする。この複体は非輪状な射影型複体であるが、
零複体にホモトピックではない。同じ複体は、零複体にホモトピックでない非輪状な
入射型複体の例にもなる。

\VerdierSemanticAnchor{ch3d-2-4-5-b}
b) $\mathcal A$ の対象の射影次元に上界がない場合でさえ、$Q$-左自由な射影型複体で、
$D(\mathcal A)$ におけるその像が $D^{-}(\mathcal A)$ に属さないものが存在しうる。
たとえば、Cartan--Eilenberg 射影複体は $Q$-左自由対象である。
より一般に、$Q$-左自由対象の任意の直和は $Q$-左自由である。

\VerdierSemanticAnchor{ch3d-2-4-6}
\medskip\noindent\textbf{2.4.6.}\quad
この節の結果を圏 $\mathcal A$ の反対圏に適用して得られる主張の定式化は読者に委ねる。
\endgroup
